% Additive Lane 4 structural proof repairs.
% This file is not included by a canonical manuscript and carries no theorem
% promotion by itself.

\section{Structural routing repairs for the quartic audit}

Throughout, \(k\) is an algebraically closed field of characteristic zero.
The statements below are candidate reader proofs for the two global routing
interfaces denoted \texttt{AUDIT-LEAD-FACT} and
\texttt{AUDIT-FOUR-LOCI}.  Their purpose is to remove shortcuts from the
argument and to make every complementary case visible.  They still require
specialist review before use in a theorem-bearing manuscript.

\subsection{Leading-image factorization via the actual normalization}

\begin{lemma}[Singular homogeneous Jacobian forces a curve image]
\label{lem:lane4-leading-image-dimension}
Let \(H=(H_1,H_2,H_3)\) be a nonzero homogeneous polynomial map of one
positive degree \(d\), and suppose
\[
 \det JH=0.
\]
Then the closure of the projective image of
\[
 [H_1:H_2:H_3]\colon\mathbf P^2\dashrightarrow\mathbf P^2
\]
has dimension at most one.
\end{lemma}

\begin{proof}
Suppose instead that the projective image has dimension two.  At a general
point \(x\) outside the base locus, the projective differential has rank two.
Choose tangent lifts \(v_1,v_2\in k^3\) such that the classes of
\(JH(x)v_1,JH(x)v_2\) are independent modulo the radial target line
\(kH(x)\).  Euler's identity gives
\[
 JH(x)x=dH(x).
\]
Consequently \(JH(x)x,JH(x)v_1,JH(x)v_2\) are linearly independent: reducing
a relation modulo \(kH(x)\) first kills the coefficients of the two tangent
vectors, and the remaining radial coefficient then vanishes.  Thus
\(\operatorname{rank}JH(x)=3\), contradicting \(\det JH=0\).
\end{proof}

\begin{proposition}[Leading-image factorization without a relative-closure shortcut]
\label{prop:lane4-leading-image-direct}
Let
\[
 H_4=(H_{4,1},H_{4,2},H_{4,3})\in k[x,y,z]_4^3
\]
be nonzero, and suppose that the closure \(C\subset \mathbf P^2\) of the
projective image of
\[
 \phi=[H_{4,1}:H_{4,2}:H_{4,3}]
 \colon \mathbf P^2\dashrightarrow \mathbf P^2
\]
is an integral curve of degree \(e\geq 1\).  Then there exist
\begin{itemize}
\item coprime homogeneous forms \(A,B\in k[x,y,z]_m\) of a common degree
  \(m\geq 1\);
\item binary forms \(h_0,h_1,h_2\in k[u,v]_e\) with no common zero such that
  \(\nu=[h_0:h_1:h_2]\colon\mathbf P^1\to C\) is the normalization map; and
\item a homogeneous form \(G\in k[x,y,z]\),
\end{itemize}
for which
\[
 H_4=G\,h(A,B),\qquad \deg G+em=4.
\]
For a nondegenerate image \(e\geq2\), the numerical possibilities are exactly
\[
 (e,m,\deg G)=(2,1,2),(2,2,0),(3,1,1),(4,1,0).
\]
The case \(e=1\) is exactly the leading-target-span-two child.
\end{proposition}

\begin{proof}
The rational map \(\phi\) dominates \(C\).  Choose a nonconstant
rational function \(f\in k(C)\setminus k\).  The pullback \(\phi^*f\)
is nonconstant on \(\mathbf P^2\), and its restriction to a sufficiently
general line remains nonconstant.  Equivalently, the restriction of
\(\phi\) to that line gives a nonconstant rational map
\(\mathbf P^1\dashrightarrow C\).  Thus \(C\) is unirational.  The
normalization \(\widetilde C\) is therefore a smooth unirational projective
curve, hence \(\widetilde C\simeq\mathbf P^1\) in characteristic zero.

Let
\[
 \nu=[h_0:h_1:h_2]\colon\mathbf P^1\longrightarrow C
\]
be the normalization map.  Because \(\nu\) is finite birational and
\(\deg C=e\), the line bundle \(\nu^*\mathcal O_C(1)\) has degree \(e\).
After choosing homogeneous coordinates \([u:v]\) on \(\mathbf P^1\), the
three sections defining \(\nu\) are therefore binary forms of degree \(e\)
without a common zero.

The rational inverse of \(\nu\) on the generic point of \(C\) gives a
rational map
\[
 \psi=\nu^{-1}\circ\phi\colon\mathbf P^2\dashrightarrow\mathbf P^1.
\]
Write \(\psi=[A:B]\) with \(A,B\) coprime homogeneous forms of one common
degree \(m\).  Since \(C\) is a curve, \(\psi\) is nonconstant and \(m\geq1\).
Put
\[
 q_i=h_i(A,B)\qquad(0\leq i\leq2).
\]
Then
\[
 [H_{4,1}:H_{4,2}:H_{4,3}]=[q_0:q_1:q_2]
\]
as rational maps.

The forms \(q_0,q_1,q_2\) have no common irreducible divisor.  Indeed, let
\(\Gamma\) be a prime divisor of \(\mathbf P^2\).  Coprimality of \(A,B\)
means that they do not both vanish at the generic point of \(\Gamma\), so
\([A:B]\) gives a point of \(\mathbf P^1(k(\Gamma))\).  If \(\Gamma\)
divided all \(q_i\), the three \(h_i\) would vanish at that point after a
field extension, contradicting basepoint-freeness of \(\nu\).

Projective equality gives a common rational multiplier
\(\lambda\in k(x,y,z)^\times\) such that \(H_{4,i}=\lambda q_i\) for all
\(i\).  Write \(\lambda=u/v\) in lowest terms.  Since every \(H_{4,i}\) is
polynomial, \(v\mid u q_i\).  Coprimality of \(u,v\) gives \(v\mid q_i\) for
all \(i\); the preceding paragraph forces \(v\in k^\times\).  Hence
\(\lambda=G\) is a polynomial.  Homogeneity makes \(G\) homogeneous, and
comparison of degrees gives
\[
 \deg G+em=4.
\]
For \(e\geq2\), enumerating the positive integral solutions yields the four
displayed triples.  If \(e=1\), the image \(C\) is a projective line, so the
three leading coordinates span a two-dimensional target subspace; conversely
a target-span-two leading map has line image unless it degenerates to the
rank-one child.
\end{proof}

\begin{corollary}[The span-three entry to the curve tree]
\label{cor:lane4-span-three-entry}
Let \(H_4\in k[x,y,z]_4^3\) satisfy \(\det JH_4=0\), and suppose its three
coordinate forms have target span three.  Then its projective image is a
nondegenerate integral curve, so
\cref{prop:lane4-leading-image-direct} applies with \(e\geq2\).
\end{corollary}

\begin{proof}
By \cref{lem:lane4-leading-image-dimension}, the image has dimension at most
one.  It cannot be a point, because then all three coordinates are
proportional, and it cannot be a line, because a line in the target imposes
one nonzero constant linear relation on the three coordinate forms.  Hence it
is a nondegenerate curve.  The image closure of an irreducible variety is
irreducible, so its reduced closure is integral.
\end{proof}

\begin{remark}[Why relative algebraic closure does not prove properness]
\label{rem:lane4-relative-closure-warning}
Relative algebraic closure is useful for detecting a primitive source
pencil, but it does not by itself prove that a parametrization of the image
curve is birational.  For example,
\[
 [x:y:z]\longmapsto[x^2:y^2:0]
\]
has image a line.  Its image field is
\(k((x/y)^2)\), while its relative algebraic closure in
\(k(x/y,z/y)\) contains \(k(x/y)\).  Rewriting the image ratios using
\(s=x/y\) produces the degree-two map \([s^2:1:0]\), not the normalization
map.  The proposition avoids this error by choosing the normalization of
\(C\) first and only then forming \(\psi=\nu^{-1}\circ\phi\).
\end{remark}

\subsection{The weighted one-variable field}

\begin{lemma}[Quartic--cubic weighted field]
\label{lem:lane4-weighted-field}
Let \(P,Q\in k[x,y,z]_4\) be algebraically independent and let
\(0\ne R\in k[x,y,z]_3\) satisfy
\[
 \operatorname{Jac}(P,Q,R)=0.
\]
Write
\[
 P=GA,\qquad Q=GB,\qquad \gcd(A,B)=1,\qquad
 n=\deg A=\deg B,
\]
and put
\[
 t=A/B,\qquad w=R^4/Q^3.
\]
Then \(w\) is algebraic over \(k(t)\).  Moreover, there are coprime
homogeneous forms \(A_0,B_0\) of a common degree \(d\) and a rational
function \(\mathcal R\in k(s)\) of degree \(e\) such that
\[
 k(t,w)=k(s),\qquad s=A_0/B_0,\qquad
 t=\mathcal R(s),\qquad n=ed.
\]
\end{lemma}

\begin{proof}
In characteristic zero, the vanishing Jacobian together with algebraic
independence of \(P,Q\) implies that \(R\) is algebraic over \(k(P,Q)\).
Choose a nonzero polynomial relation and retain one source-homogeneous
component.  Every monomial \(P^iQ^jR^m\) in that component satisfies
\[
 4(i+j)+3m=M.
\]
Thus all exponents \(m\) are congruent modulo four.  After removing the
common residual power of \(R\), the relation has the form
\[
 \sum_{\ell} C_\ell(P,Q)R^{4\ell}=0,
 \qquad \deg C_\ell=N-3\ell.
\]
Dividing by \(Q^N\) gives a polynomial relation for
\(w=R^4/Q^3\) over \(k(t)\).

The field \(k(t,w)\) has transcendence degree one and is the function field
of a curve dominated by \(\mathbf P^2\).  Pull back a nonconstant rational
function on that curve and restrict it, together with the dominant rational
map, to a sufficiently general line.  The resulting nonconstant map from
\(\mathbf P^1\) makes the curve unirational; hence it is rational.  Therefore
\(k(t,w)=k(s)\).  Represent \(s\) by a reduced homogeneous pencil
\(A_0/B_0\) of common degree \(d\).

Write \(\mathcal R=U/V\), with coprime binary forms \(U,V\) of degree \(e\).
The substitutions \(U(A_0,B_0)\) and \(V(A_0,B_0)\) remain coprime: the
binary resultant identities express powers of \(A_0\) and \(B_0\) in their
ideal, so a common prime divisor would divide both \(A_0\) and \(B_0\).
Since
\[
 A/B=U(A_0,B_0)/V(A_0,B_0)
\]
and both ratios are reduced, comparison of degrees gives \(n=ed\).
\end{proof}

\subsection{Four structural loci with the valuation complements retained}

\begin{proposition}[Four-loci cover]
\label{prop:lane4-four-loci}
Under the hypotheses of \cref{lem:lane4-weighted-field}, every leading
triple \((P,Q,R)\) lies in at least one of the following loci:
\begin{enumerate}
\item \(P,Q,R\) belong to a polynomial ring in two independent linear forms;
\item \(G=1\) and
\[
 P=U(A_0,B_0),\qquad Q=V(A_0,B_0),
\]
where \(A_0,B_0\) are coprime quadrics and \(U,V\) are binary quadratics;
\item \(G=1\), the reduced pencil \(A/B\) is composition-primitive, and the
pencil contains a fourth power \(\ell^4\);
\item the reduced pencil \(A/B\) is composition-primitive, \(G\ne1\), and
every irreducible component of \(G\) is supported on a special fiber of
\(A/B\).
\end{enumerate}
The assertion is a cover, not a claim that the four loci are disjoint.
\end{proposition}

\begin{proof}
Use the notation \(n=ed\) from \cref{lem:lane4-weighted-field}.

Suppose first that \(e>1\).  Since \(n=4-\deg G\), the possibilities with
\(d\geq1\) are
\[
\begin{array}{c|c|c}
\deg G&n&(e,d)\\ \hline
0&4&(4,1),(2,2)\\
1&3&(3,1)\\
2&2&(2,1).
\end{array}
\]
In every \(d=1\) case, \(s=A_0/B_0\) is a pencil of two independent linear
forms, and the reduced forms \(A,B\) are binary in those forms.  Let
\(\Gamma\) occur in \(G\) with multiplicity \(\mu\).  If \(s\) is
nonconstant on \(\Gamma\), then \(B\) and every nonzero rational function of
\(s\) are units at the generic point of \(\Gamma\).  Taking
\(\Gamma\)-valuations in
\[
 R^4/Q^3=\sigma(s)
\]
gives
\[
 4\nu_\Gamma(R)=3\mu.
\]
For the displayed \(d=1\) cases with \(G\ne1\), one has
\(1\leq\mu\leq\deg G\leq2\), so this is impossible.  Hence every component of
\(G\) is a fiber line of \(s\), and \(G,P,Q\) are binary in the same two
linear forms.  The identity \(R^4=Q^3\sigma(s)\) shows that \(R^4\) belongs
to their rational function field.  Choose a third independent linear form
\(c\).  Writing \(R\) as a polynomial in \(c\), the leading
\(c\)-coefficient of \(R^4\) is the fourth power of the leading
coefficient of \(R\); it cannot vanish in a domain.  Since the right-hand
side has \(c\)-degree zero, \(R\) has \(c\)-degree zero and is binary.  This is locus (1).  The sole nonbinary composite case is
\((e,d)=(2,2)\) with \(G=1\), which is locus (2).

Now suppose \(e=1\), so \(w=\rho(t)\).  First take \(G=1\).  For a finite
fiber \(A-\xi B\), set \(c_\xi=\operatorname{ord}_\xi\rho\); at infinity set
\[
 c_\infty=\operatorname{ord}_\infty\rho+3.
\]
If an irreducible component \(\Gamma\) occurs in the corresponding quartic
fiber with multiplicity \(m\), valuation gives
\[
 4\nu_\Gamma(R)=c_\xi m.
\]
Consequently every \(c_\xi\geq0\).  The degree-zero divisor identity for
\(\rho\), with the denominator adjustment at infinity, gives
\[
 \sum_{\xi\in\mathbf P^1}c_\xi=3.
\]
At least one \(c_\xi\) is odd.  For that fiber,
\(4\mid c_\xi m\), hence \(4\mid m\) for every irreducible component.
Because the total degree is four, the fiber is \(\ell^4\).  This is
locus (3).

Finally take \(e=1\) and \(G\ne1\).  Let \(\Gamma\mid G\) have multiplicity
\(\mu\).  If \(t=A/B\) is nonconstant on \(\Gamma\), then \(A,B\), and
\(\rho(t)\) are units at its generic point, and valuation gives
\[
 4\nu_\Gamma(R)=3\mu.
\]
But \(1\leq\mu\leq\deg G\leq3\), which is impossible.  Therefore \(t\) is
constant on every component of \(G\), so each such component is supported on
a special fiber.  This is locus (4).
\end{proof}

\begin{remark}[Ownership and signs]
\label{rem:lane4-four-loci-ownership}
For a disjoint case tree, first extract the full gcd \(G\), then test whether
the resulting field generator is binary, then assign the quadratic-source
and fourth-power cases, and only afterwards call the fixed-component owner.
The nonnegativity \(c_\xi\geq0\) was used only in the coprime case \(G=1\).
For \(G\ne1\), poles of \(\rho\) may be cancelled by \(G^3\), so imposing
\(c_\xi\geq0\) there would incorrectly delete fixed-component branches.
\end{remark}
